\documentclass[12pt,a4paper]{article}
\usepackage[utf8]{inputenc}
\usepackage{amsmath, amssymb, amsfonts}
\usepackage{geometry}
\geometry{margin=2.5cm}
\usepackage{hyperref}
\usepackage{pgfplots}
\usepackage{fancyhdr}
\usepackage{listing}
\usepackage{listings}
\usepackage{titling}
\usepackage{float}
\usepackage{titlesec}

\titleformat{\section}
  {\normalfont\Large\bfseries\centering}  % Format: font size, bold, centered
  {\thesection}{1em}{}
\pagestyle{fancy}
\pgfplotsset{compat=1.18}
\fancyhead{}
\fancyfoot{}
\title{Conoscersi per valorizzarsi}
\pretitle{\begin{center}\LARGE}
\posttitle{\par\large Career Labs\end{center}}
\author{}
\date{}
\fancyhead[L]{Gabriele Perna}          % Left side
\fancyhead[C]{Conoscersi per valorizzarsi}     % Center
\fancyhead[R]{2026/2027}             % Right side
\fancyfoot[C]{\thepage}           % Page number in the center of footer
\let\oldsection\section
\renewcommand{\section}{\clearpage\oldsection}

\begin{document}

\maketitle
\tableofcontents

\section{Informazioni e risorse}
\textbf{Nominativo consulente:} Giulia Verdigi.\\
\textbf{Email consulente:} giulia.verdigi@ospiti.unipi.it\\
\textbf{Instagram consulente:} gverdigi\_psicologa\\\\
\textbf{Consulenza} \textbf{online} \textbf{(gratuita) sui colloqui di lavoro:}
\href{https://alumni.unipi.it/mentoring}{https://alumni.unipi.it/mentoring}

\section{Come mai questo laboratorio?}
Molte persone non vedono il proprio valore e non si valorizzano.\\
Tanti soggetti non esprimono i propri hobby e valori nascosti quando invece dovrebbero farlo, in quanto li rende unici e li aiuta a distinguersi.\\
Un recruiter potrebbe ricordarsi meglio di noi proprio per via di alcune informazioni uniche a noi.\\\\
Inoltre, è necessario essere consapevoli di quello che siamo e dei nostri interessi in modo da seguire una rotta adatta a noi e garantirci un benessere nel mondo dello studio e del lavoro.\\\\
Altri esempi di vantaggi dell'autoconsapevolezza sono:
\begin{itemize}
  \item Conoscere i propri limiti
  \item Empatia
  \item Cogliere le opportunità
  \item Decisioni consapevoli
  \item Adattabilità
  \item Migliori prestazioni
\end{itemize}

\subsection{Cosa è il talento?}
Il talento è un abilità, inclinazione naturale.\\
Ogni persona ha un talento, tocca alla stessa prendersene cura, bisogna però esserne consapevole.\\
Un talento corrisponde a delle skills, per questo bisogna conoscerli ed esporli.\\
Per esempio, il semplice talento nel gioco delle freccette può essere tramutato come:
\begin{itemize}
  \item Precisione
  \item Equilibrio
  \item Stabilità
  \item Concentrazione
\end{itemize}
Queste sono tutte skill assolutamente valide per il mondo del lavoro.

\section{Combattere i pensieri negativi}
Se ci riteniamo un fallimento o abbiamo la \textbf{sindrome dell'impostore}, per combattere questi pensieri negativi bisogna tenere sempre a mente alcune cose importanti:\
\begin{itemize}
  \item \textbf{I traguardi gia raggiunti}: se siamo riusciti a fare A, che già è impressionante, allora cosa ci vieta di fare B?
  \item \textbf{Autostima}: non importa cosa accadrà, mi vorrò sempre bene.
\end{itemize}

\section{Intelligenza emotiva}
L'intelligenza emotiva si ottiene effettuando questi passaggi:
\begin{enumerate}
  \item Autocoscienza
  \item Gestione delle emozioni
  \item Motivare se stessi
  \item Riconoscere le emozioni negli altri (empatia)
  \item Gestione delle relazioni
\end{enumerate}
Chi possiede questa intelligenza ha una maggiore \textbf{attenzione ecologica}, maggiore \textbf{empatia}, \textbf{successo}, \textbf{team performance}, meno \textbf{stress} e una migliore vita familiare e di coppia.\\
Tali persone hanno una migliore \textbf{work-family life}.

\section{Esercizio: descriviti con 3 aggettivi}
NON DIRE ASSOLUTAMENTE:
\begin{itemize}
  \item Solare
  \item Disponibile
  \item Dinamico
\end{itemize}
Questi sono banali e over-used.\\\\
Non necessariamente devono essere tutti positivi, potrebbero essere anche legati a difetti.\\
Nell'esporre i difetti, bisogna assicurarsi che non ci rovinino l'opportunità, bisogna esporre in che modo stiamo lavorando per risolverla o attenuarla (o come la usiamo a nostro favore).

\section{Matrice SWOT e domande personali}
Per conoscersi meglio, possiamo lavorare sulla \textbf{matrice SWOT}, la quale permette di analizzarci sotto diversi aspetti, quali:
\begin{itemize}
  \item \textbf{Strengths:} punti di forza, caratteristiche personali o risorse che danno un vantaggio
  \item \textbf{Weaknesses:} punti di debolezza, caratteristiche personali in cui c'è margine di miglioramento
  \item \textbf{Opportunities:} opportunità, fattori esterni che possono essere usati a nostro favore
  \item \textbf{Threats:} minacce, fattori esterni che possono rappresentare ostacoli
\end{itemize}
Ci sono molte domande da porsi per aumentare la nostra consapevolezza e riuscire a presentarsi meglio, vediamole.
\subsection{Valori}
\begin{itemize}
  \item Quali sono le cose più importanti per te nella vita?
  \item Cosa ti motiva a fare le scelte che fai?
  \item Quali sono le tue priorità?
  \item A cosa non puoi rinunciare?
\end{itemize}
\subsection{Interessi}
\begin{itemize}
  \item Quali attività ti fanno sentire più vivo ed energetico?
  \item In cosa spendi il tuo tempo libero?
  \item Quali argomenti o tematiche trovi più interessanti da esplorare?
  \item Cosa ti piacerebbe imparare o approfondire di più?
\end{itemize}
\subsection{}
\end{document}